% Part: incompleteness
% Chapter: arithmetization-syntax
% Section: proofs-in-lk

\documentclass[../../../include/open-logic-section]{subfiles}

\begin{document}

\olfileid[fr]{inc}{art}{plk}
\olsection{\usetoken{P}{derivation} dans $\Log{LK}$}

\begin{explain}
Pour arithmétiser les !!{derivation}s, nous devons les représenter
par des nombres. Une !!{derivation} est un arbre de séquents dont
chaque inférence porte aussi une étiquette. Une représentation
récursive s'impose donc : nous représentons une !!{derivation} par
un uplet dont les composantes sont le séquent final, l'étiquette
et les représentations des sous-!!{derivation}s qui aboutissent aux
prémisses de la dernière inférence.
\end{explain}

\begin{defn}
Si $\Gamma$ est une suite finie d'!!{sentence}s,
$\Gamma = \tuple{!A_1, \dots, !A_n}$, alors
$\Gn{\Gamma} = \tuple{\Gn{!A_1}, \dots, \Gn{!A_n}}$.

Si $\Gamma \Sequent \Delta$ est un séquent, son numéro de G\"odel est
\[
\Gn{\Gamma \Sequent \Delta} = \tuple{\Gn{\Gamma}, \Gn{\Delta}}
\]

Si $\pi$ est !!a{derivation} dans $\Log{LK}$, son numéro de
G\"odel $\Gn{\pi}$ est défini comme suit :
\begin{enumerate}
\item Si $\pi$ se réduit au séquent initial
  $\Gamma \Sequent \Delta$, alors $\Gn{\pi}$ vaut
  \[
    \tuple{0, \Gn{\Gamma \Sequent \Delta}}.
  \]
\item Si $\pi$ se termine par une inférence à une ou deux prémisses,
  de conclusion $\Gamma \Sequent \Delta$, et si $\pi_1$ et $\pi_2$
  désignent les sous-dérivations immédiates aboutissant aux prémisses
  de cette dernière inférence, alors $\Gn{\pi}$ vaut respectivement
  \begin{align*}
    & \tuple{1, \Gn{\pi_1}, \Gn{\Gamma \Sequent \Delta}, k} \text{ ou}\\
    & \tuple{2, \Gn{\pi_1}, \Gn{\pi_2}, \Gn{\Gamma \Sequent \Delta}, k},
  \end{align*}
  où $k$ est déterminé, selon la règle employée dans la dernière
  inférence, par le tableau suivant :

  \begin{tabular}{lccccccc}
    \text{Règle :} & \LeftR{\Weakening} & \RightR{\Weakening} &
    \LeftR{\Contraction} & \RightR{\Contraction} &
    \LeftR{\Exchange} & \RightR{\Exchange} \\
    $k$: & 1 & 2 & 3 & 4 & 5 & 6 \\[2ex]
    \text{Règle :} & \LeftR{\lnot} & \RightR{\lnot} &
    \LeftR{\land} & \RightR{\land} &
    \LeftR{\lor} & \RightR{\lor} \\
    $k$: & 7 & 8 & 9 & 10 & 11 & 12 \\[2ex]
    \text{Règle :} & \LeftR{\lif} & \RightR{\lif} &
    \LeftR{\lforall} & \RightR{\lforall} &
    \LeftR{\lexists} & \RightR{\lexists} \\
    $k$: & 13 & 14 & 15 & 16 & 17 & 18 \\[2ex]
    \text{Règle :} & \Cut & = \\
    $k$: & 19 & 20
  \end{tabular}
\end{enumerate}
\end{defn}

\begin{ex}
  Considérons la !!{derivation} très simple
  \begin{prooftree}
    \Axiom$!A \fCenter !A$
    \RightLabel{\LeftR{\land}}
    \UnaryInf$!A \land !B \fCenter !A$
    \RightLabel{\RightR{\lif}}
    \UnaryInf$\fCenter (!A \land !B) \lif !A$
  \end{prooftree}
  Le numéro de G\"odel de la dérivation réduite au séquent initial
  est $p_0 = \tuple{0, \Gn{!A \Sequent !A}}$. Le numéro de G\"odel
  de la !!{derivation} dont la conclusion résulte de
  $\LeftR{\land}$ est
  $p_1 = \tuple{1, p_0, \Gn{!A \land !B \Sequent !A}, 9}$ : le
  $1$ indique que $\LeftR{\land}$ a une seule prémisse; viennent
  ensuite le numéro de G\"odel de la conclusion
  $!A \land !B \Sequent !A$ et le $9$ qui code $\LeftR{\land}$.
  Le numéro de G\"odel de la !!{derivation} entière est donc
  $\tuple{1, p_1, \Gn{\Sequent (!A \land !B) \lif !A}, 14}$,
  c'est-à-dire
  \[
  \tuple{1, \tuple{1, \tuple{0, \Gn{!A \Sequent !A}}, \Gn{!A \land !B \Sequent !A}, 9},
    \Gn{\Sequent (!A \land !B) \lif !A}, 14}.
  \]
\end{ex}

\begin{explain}
Une fois cette représentation des !!{derivation}s fixée, nous devons
montrer que nous pouvons les manipuler par des opérations récursives
primitives et exprimer de la même façon leurs propriétés et relations
essentielles. Certaines opérations sont simples : si $p$ est le
numéro de G\"odel d'!!a{derivation},
$\fn{EndSequent}(p) = (p)_{(p)_0+1}$ donne le numéro de G\"odel de
son séquent final et $\fn{LastRule}(p) = (p)_{(p)_0+2}$ le code de sa
dernière règle. La propriété $\fn{Sequent}(s)$, définie par
\begin{multline*}
  \fn{SeqCode}(s)\land\len{s}=2\land\fn{SeqCode}((s)_0)\land\fn{SeqCode}((s)_1)\land{}\\ \bforall{i<\len{(s)_0} + \len{(s)_1}}{\fn{Sent}(((s)_0 \concat (s)_1)_i)}
\end{multline*}
est vraie de $s$ si et seulement si $s$ est le numéro de G\"odel
d'un séquent composé d'!!{sentence}s. D'autres propriétés sont bien
plus difficiles à traiter; nous en esquisserons au moins la méthode.
Notre objectif est de montrer que la relation « $\pi$ est
!!a{derivation} de~$!A$ à partir de~$\Gamma$ » est récursive
primitive sur les numéros de G\"odel de $\pi$ et de~$!A$.
\end{explain}

\begin{prop}\ollabel{prop:followsby}
  La propriété $\fn{Correct}(p)$, vraie si et seulement si la
  dernière inférence de la !!{derivation}~$\pi$ de numéro de
  G\"odel~$p$ est correcte, est récursive primitive.
\end{prop}

\begin{proof}
  $\Gamma \Sequent \Delta$ est un séquent initial s'il existe soit
  !!a{sentence}~$!A$ tel que ce séquent soit $!A \Sequent !A$,
  soit un terme clos~$t$ tel qu'il soit $\emptyset \Sequent \eq[t][t]$.
  En termes de numéros de G\"odel, $\fn{InitSeq}(s)$ est vraie si et
  seulement si
  \begin{align*}
  \bexists{x < s}{} (\fn{Sent}(x) & \land
  s = \tuple{\tuple{x},\tuple{x}})
  \lor {}\\
  \bexists{t<s}{} (\fn{ClTerm}(t) & \land
  s = \tuple{0, \tuple{\Gn{{\eq}(} \concat t \concat \Gn{,} \concat t \concat \Gn{)}}}).
  \end{align*}

  Il faut encore montrer que, pour chaque règle d'inférence~$R$,
  la relation $\fn{FollowsBy}_R(p)$ est récursive primitive. Elle
  est vraie si et seulement si $p$ est le numéro de G\"odel
  d'!!a{derivation}~$\pi$ dont le séquent final résulte d'une
  application correcte de~$R$ aux sous-!!{derivation}s immédiates
  de~$\pi$.

  Le cas de la règle \RightR{\land} est simple. Si $\pi$ se termine
  par une application correcte de cette règle, elle a la forme
  \begin{prooftree}
    \AxiomC{}
    \RightLabel{$\pi_1$}
    \Deduce$\Gamma \fCenter \Delta, !A$

    \AxiomC{}
    \RightLabel{$\pi_2$}
    \Deduce$\Gamma \fCenter \Delta, !B$

    \RightLabel{\RightR\land}
    \BinaryInf$\Gamma \fCenter \Delta, !A \land !B$
  \end{prooftree}
  La dernière inférence de la !!{derivation}~$\pi$ est donc une
  application correcte de $\RightR{\land}$ s'il existe des suites
  d'!!{sentence}s $\Gamma$ et $\Delta$ et deux !!{sentence}s $!A$
  et~$!B$ tels que le séquent final de $\pi_1$ soit
  $\Gamma \Sequent \Delta, !A$, celui de $\pi_2$ soit
  $\Gamma \Sequent \Delta, !B$, et celui de $\pi$ soit
  $\Gamma \Sequent \Delta, !A \land !B$. Il reste à traduire ces
  conditions en numéros de G\"odel. Si
  $s = \Gn{\Gamma \Sequent \Delta}$, alors
  $(s)_0 = \Gn{\Gamma}$ et $(s)_1 = \Gn{\Delta}$. Par conséquent,
  $\fn{FollowsBy}_{\RightR{\land}}(p)$ est vraie si et seulement si
\begin{align*}
& \bexists{g < p}{\bexists{d < p}{\bexists{a < p}{\bexists{b < p}{\quad}}}} \\
& \qquad \fn{EndSequent}(p) =
  \tuple{g, d \concat
    \tuple{\Gn{(} \concat a \concat \Gn{\land} \concat b \concat \Gn{)}}} \land {} \\
& \qquad \fn{EndSequent}((p)_1) = \tuple{g, d \concat \tuple{a}} \land {}\\
& \qquad \fn{EndSequent}((p)_2) = \tuple{g, d \concat \tuple{b}} \land {}\\
& \qquad (p)_0 = 2 \land \fn{LastRule}(p) = 10.
\end{align*}
Les lignes successives expriment l'existence d'une suite~$\Gamma$
de numéro de G\"odel~$g$, d'une suite~$\Delta$ de numéro~$d$,
d'!!a{formula}~$!A$ de numéro~$a$ et d'!!a{formula}~$!B$ de
numéro~$b$; puis que le séquent final de $\pi$ est
$\Gamma \Sequent \Delta, !A \land !B$, celui de $\pi_1$ est
$\Gamma \Sequent \Delta, !A$, celui de $\pi_2$ est
$\Gamma \Sequent \Delta, !B$; enfin, que $\pi$ possède deux
sous-dérivations immédiates et que sa dernière règle est
$\RightR\land$, codée par~$10$.

La dernière inférence de~$\pi$ est une application correcte de
$\RightR\lexists$ s'il existe des suites $\Gamma$ et $\Delta$,
!!a{formula}~$!A$, une variable~$x$ et un terme~$t$, tels que
le séquent final de $\pi$ soit
$\Gamma \Sequent \Delta, \lexists[x][!A]$ et celui de $\pi_1$
soit $\Gamma \Sequent \Delta, \Subst{!A}{t}{x}$. En numéros de
G\"odel, $\fn{FollowsBy}_{\RightR\lexists}(p)$ est donc vraie si et
seulement si
\begin{align*}
  & \bexists{g<p}{\bexists{d<p}{\bexists{a<p}{\bexists{x<p}{\bexists{t<p}{\quad}}}}}\\
  & \qquad \fn{EndSequent}(p) =
  \tuple{
    g,
    d \concat \tuple{\Gn{\lexists} \concat x \concat a}
  } \land {}\\
  & \qquad \fn{EndSequent}((p)_1) =
  \tuple{
    g,
    d \concat \tuple{\fn{Subst}(a, t, x)}
  } \land {}\\
  & \qquad \fn{Frm}(a)\land\fn{Var}(x)\land\fn{ClTerm}(t)\land{}\\ & \qquad (p)_0 = 1 \land \fn{LastRule}(p) = 18.
\end{align*}
Appelons $\fn{ShapeLK}(p)$ la condition $\fn{SeqCode}(p)$ et, selon $(p)_0=0,1,2$, respectivement $\len{p}=2,4,5$ ; toute autre arité est exclue. Nous définissons alors $\fn{Correct}(p)$ par
\begin{multline*}
  \fn{ShapeLK}(p)\land\fn{Sequent}(\fn{EndSequent}(p)) \land {}\\
  [(\fn{LastRule}(p) = 1 \land
  \fn{FollowsBy}_{\LeftR\Weakening}(p)) \lor \dots \lor {}\\
  (\fn{LastRule}(p) = 20 \land \fn{FollowsBy}_{\eq}(p)) \lor {}\\
  (p)_0 = 0 \land \fn{InitSeq}(\fn{EndSequent}(p))]
\end{multline*}
La première ligne garantit que le séquent final de la dérivation
codée par~$p$ est bien composé d'!!{sentence}s. La dernière traite
le cas où $p$ code uniquement un séquent initial.
\end{proof}

\begin{prob}
  Définir les propriétés suivantes, comme dans la
  \olref[inc][art][plk]{prop:followsby} :
  \begin{enumerate}
  \item $\fn{FollowsBy}_{\Cut}(p)$;
  \item $\fn{FollowsBy}_{\LeftR\lif}(p)$;
  \item $\fn{FollowsBy}_{\eq}(p)$;
  \item $\fn{FollowsBy}_{\RightR{\lforall}}(p)$.
  \end{enumerate}
  Pour la dernière, il faudra aussi montrer qu'on peut vérifier de
  façon récursive primitive si la dernière inférence de la
  !!{derivation} de numéro de G\"odel~$p$ satisfait la condition sur
  la variable propre : la constante propre~$a$ de
  $\RightR{\lforall}$ ne doit pas apparaître dans le séquent final.
\end{prob}

\begin{prop}
  \ollabel{prop:deriv}
  La relation $\fn{Deriv}(p)$, vraie si et seulement si $p$ est le
  numéro de G\"odel d'!!a{derivation} correcte~$\pi$, est
  récursive primitive.
\end{prop}

\begin{proof}
  !!^a{derivation}~$\pi$ est correcte si chacune de ses inférences
  est une application correcte d'une règle, autrement dit si chacune
  de ses sous-!!{derivation}s se termine par une inférence correcte.
  Ainsi, $\fn{Deriv}(p)$ est vraie si et seulement si
  \[
  \bforall{i<\len{\fn{SubtreeSeq}(p)}}{\fn{Correct}((\fn{SubtreeSeq}(p))_i)}.
  \]
\end{proof}

\begin{prop}
Supposons que $\Gamma$ soit un ensemble récursif primitif
d'!!{sentence}s. Alors la relation $\Prf[\Gamma](x, y)$, exprimant
que « $x$ est le code d'!!a{derivation}~$\pi$ de
$\Gamma_0 \Sequent !A$ pour un certain ensemble fini
$\Gamma_0 \subseteq \Gamma$, et $y$ est le numéro de G\"odel de~$!A$ »,
est récursive primitive.
\end{prop}

\begin{proof}
Supposons que l'appartenance à~$\Gamma$ soit donnée, sur les codes,
par le prédicat récursif primitif~$R_\Gamma(y)$. Nous devons montrer
que $\Prf[\Gamma](x, y)$, vraie si et seulement si $y$ est le
numéro de G\"odel d'un énoncé~$!A$ et que $x$ est le code d'!!a{derivation}
dans $\Log{LK}$ dont le séquent final est
$\Gamma_0 \Sequent !A$, est récursive primitive.

D'après la proposition précédente, la propriété $\fn{Deriv}(x)$,
qui dit que $x$ est le code d'!!a{derivation} correcte~$\pi$ dans
$\Log{LK}$, est récursive primitive. Si $x$ est un tel code,
$\fn{EndSequent}(x)$ code le séquent final de~$\pi$; ainsi,
$(\fn{EndSequent}(x))_0$ code son côté gauche et
$(\fn{EndSequent}(x))_1$ son côté droit. Nous pouvons donc exprimer
« le côté droit du séquent final de~$\pi$ est~$!A$ » par
$\len{(\fn{EndSequent}(x))_1} = 1 \land ((\fn{EndSequent}(x))_1)_0 =
y$. Le côté gauche est nécessairement fini; il reste seulement à
exprimer que chacun de ses énoncés appartient à~$\Gamma$. Nous
pouvons ainsi définir $\Prf[\Gamma](x, y)$ par
\begin{align*}
\Prf[\Gamma](x, y) \defiff {}&
\fn{Deriv}(x) \land {} \\
& \bforall{i <
  \len{(\fn{EndSequent}(x))_0}}{R_\Gamma(((\fn{EndSequent}(x))_0)_i)} \land {}\\
& \len{(\fn{EndSequent}(x))_1} = 1 \land ((\fn{EndSequent}(x))_1)_0 = y.
\end{align*}
\end{proof}

\begin{explain}
Pour vérifier un code de dérivation, on contrôle d'abord la forme du
tuple à chaque nœud, puis les deux suites d'énoncés qui composent
son séquent. À un nœud initial, on vérifie le séquent initial.
Aux autres nœuds, on vérifie la règle indiquée par le dernier
champ du tuple à partir des séquents des prémisses. Les schémas
qui emploient un terme clos exigent le test $\fn{ClTerm}$ sur
son code. Enfin, pour $\Prf[\Gamma](x,y)$, le séquent final doit
avoir pour seul énoncé à droite celui de code~$y$ et pour
énoncés à gauche des éléments de~$\Gamma$.
\end{explain}

\end{document}
